\begin{textsample}{POS Dim 1 – ap – Score 45.00 – ap/ap\_ef\_17.txt}  \label{ex:f1_pos_168}
4.3.4.1.
História “ História é vital para a formação da cidadania porque nos mostra que para compreender o que está acontecendo no presente é preciso entender quais foram os \textbf{caminhos} percorridos pela sociedade. ” ( Boris Fausto - Historiador ).
Como forma específica de produção de conhecimentos, a História tem \textbf{relevante} importância para se alcançar os objetivos propostos pela educação voltada para emancipação, bem como, compreender as semelhanças e as diferenças e as transformações no modo de vida social, cultural e \textbf{econômico} de sua \textbf{localidade}, ou seja, a garantia de formação de um educando com uma concepção de cidadania, e sabedor do seu papel enquanto sujeito social.
Ao fazer parte das etapas de Alfabetização e \textbf{letramento} o ensino de História, proporciona que o \textbf{aluno} construa e desenvolva um perfil \textbf{inerente} às \textbf{demandas} de uma formação necessária a enfrentar desafios e dar soluções, transformando -se em sujeito de sua existência.
Essa autonomia intelectual subsidiará o educando a desenvolver formas de \textbf{compreensão} e maior reflexão sobre a realidade, bem como, a expressão de ideias e elaboração de conceitos, e de atuação, enquanto cidadãos conscientes de direitos.
Para alcançar os objetivos, a organização e estruturação do ensino de História, deve considerar o contexto vivido e vivenciado pela criança em seu cotidiano.
Sendo considerado para tal intento um planejamento que possibilite ao \textbf{aluno} a observação do seu contexto \textbf{atual} relacionando -o a fatos de seu passado.
O estímulo à pesquisa e a investigação de diferentes sujeitos e processos históricos deve ser considerado importante, assim como, \textbf{debates} das diferentes relações dos mais diversos agrupamentos sociais em diferentes épocas e lugares.
No Ensino Fundamental de nove anos, o conhecimento deve refletir o aspecto significativo dos conceitos e \textbf{categorias} básicas desse campo : tempo, espaço, documento, sujeito histórico, fato histórico.
Ao se deparar com diversas situações do dia a dia o educando vai construindo e ressignificando seus conceitos previamente elaborados.
De acordo com a BNCC As questões que nos levam a pensar a História como um saber necessário para a formação das crianças e jovens na escola são as originárias do tempo presente.
O passado que deve impulsionar a dinâmica do ensino-aprendizagem no Ensino Fundamental é aquele que dialoga com o tempo \textbf{atual}.
A relação passado/presente não se \textbf{processa} de forma automática, pois \textbf{exige} o conhecimento de referências \textbf{teóricas} capazes de trazer inteligibilidade aos objetos históricos selecionados.
Um objeto só se torna documento quando apropriado por um narrador que a ele confere sentido, tornando -o capaz de expressar a dinâmica da vida das sociedades.
Portanto, o que nos interessa no conhecimento histórico é perceber a forma como os indivíduos construíram, com diferentes linguagens, suas narrações sobre o mundo em que viveram e vivem, suas instituições e organizações sociais.
A BNCC de História no Ensino Fundamental – Anos Iniciais contempla, antes de mais nada, a construção do sujeito.
O processo tem início quando a criança toma consciência da existência de um “ Eu ” e de um “ Outro ”.
O exercício de separação dos sujeitos é um \textbf{método} de conhecimento, uma \textbf{maneira} pela qual o indivíduo toma consciência de si, desenvolvendo a \textbf{capacidade} de administrar a sua vontade de \textbf{maneira} autônoma, como parte de uma família, uma comunidade e um corpo social.
Esse processo de constituição do sujeito é longo e \textbf{complexo}.
Os indivíduos desenvolvem sua percepção de si e do outro em meio a vivências cotidianas, identificando o seu lugar na família, na escola e no espaço em que vivem.
O aprendizado, ao longo do Ensino Fundamental – Anos Iniciais, torna -se mais \textbf{complexo} à medida que o sujeito reconhece que existe um “ Outro ” e que cada um apreende o mundo de forma particular.
A percepção da distância entre objeto e pensamento é um passo necessário para a autonomia do sujeito, tomado como produtor de diferentes linguagens.
É ela que funda a relação do sujeito com a sociedade.
Nesse sentido, a História \textbf{depende} das linguagens com as quais os seres humanos se comunicam, entram em conflito e negociam.
A busca de autonomia também \textbf{exige} reconhecimento das bases da \textbf{epistemologia} da História, a saber : a natureza compartilhada do sujeito e do objeto de conhecimento, o conceito de tempo histórico em seus diferentes ritmos e durações, a concepção de documento como suporte das relações sociais, as várias linguagens por meio das quais o ser humano se apropria do mundo.
Enfim, percepções capazes de responder aos desafios da prática historiadora presente dentro e fora da sala de \textbf{aula}.
Todas essas considerações de ordem \textbf{teórica} devem considerar a experiência dos \textbf{alunos} e professores, tendo em \textbf{vista} a realidade social e o universo da comunidade escolar, bem como seus referenciais históricos, sociais e culturais.
Ao promover a diversidade de \textbf{análises} e proposições, \textbf{espera} -se que os \textbf{alunos} construam as próprias interpretações, de forma fundamentada e rigorosa. \textbf{Convém} destacar as temáticas voltadas para a diversidade cultural e para as múltiplas configurações \textbf{identitárias}, destacando -se as \textbf{abordagens} relacionadas à história dos povos indígenas originários e africanos.
Ressalta -se, também, na formação da sociedade brasileira, a presença de diferentes povos e culturas, suas contradições sociais e culturais e suas articulações com outros povos e sociedades.
A inclusão dos temas obrigatórios definidos pela legislação vigente, tais como a história da África e das culturas afro-brasileira e indígena, deve ultrapassar a dimensão puramente retórica e permitir que se defenda o estudo dessas populações como artífices da própria história do Brasil.
A relevância da história desses grupos humanos reside na possibilidade de os estudantes compreenderem o papel das alteridades presentes na sociedade brasileira, comprometerem -se com elas e, ainda, perceberem que existem outros referenciais de produção, circulação e transmissão de conhecimentos, que podem se entrecruzar com aqueles considerados consagrados nos espaços formais de produção de saber.
Problematizando a ideia de um “ Outro ”, \textbf{convém} observar a presença de uma percepção estereotipada naturalizada de diferença, ao se tratar de indígenas e africanos.
Essa problemática está associada à produção de uma história brasileira marcada pela imagem de nação constituída nos moldes da colonização europeia.
Por todas as \textbf{razões} apresentadas, \textbf{espera} -se que o conhecimento histórico seja tratado como uma forma de pensar, entre várias ; uma forma de indagar sobre as coisas do passado e do presente, de construir explicações, \textbf{desvendar} significados, compor e decompor interpretações, em movimento contínuo ao longo do tempo e do espaço.
Enfim, trata -se de transformar a história em ferramenta a serviço de um discernimento maior sobre as experiências humanas e as sociedades em que se vive.
Retornando ao ambiente escolar, a BNCC pretende estimular ações nas quais professores e \textbf{alunos} sejam sujeitos do processo de ensino e aprendizagem.
Nesse sentido, eles próprios devem assumir uma atitude historiadora diante dos conteúdos propostos no âmbito do Ensino Fundamental.
Cumpre destacar que os critérios de organização das habilidades na BNCC ( com a explicitação dos objetos de conhecimento aos quais se relacionam e do agrupamento desses objetos em unidades temáticas ) expressam um arranjo possível ( dentre outros ).
Portanto, os agrupamentos propostos não devem ser tomados como modelo obrigatório para o desenho dos currículos.
Nos anos iniciais, o ensino deste componente curricular propõe uma formação por \textbf{competências} e habilidades, tendo por objetivos de acordo com os \textbf{Parâmetros} Curriculares Nacionais, possibilitar ao \textbf{aluno} : I – Identificar o próprio grupo de convívio e as relações que se estabelecem em outros tempos e espaços ; II – Organizar repertórios que lhes permitam localizar acontecimentos numa multiplicidade de tempo, de modo a formular explicações para questões do presente e do passado ; III – Conhecer e respeitar modos de vida de diferentes grupos sociais, em diversos tempos e espaços, em suas manifestações culturais, \textbf{econômicas}, políticas e sociais, reconhecendo semelhanças e diferenças entre eles ; IV – Reconhecer mudanças e permanência nas vivencias humanas, presentes na realidade em outras comunidades, próximas ou distantes no tempo e no espaço ; V – Questionar sua realidade identificando alguns \textbf{problemas}, refletindo sobre algumas possíveis soluções, reconhecendo formas de atuação política institucional e organizações coletivas da sociedade civil ; VI – Utilizar \textbf{métodos} de pesquisa e produção de textos de conteúdos históricos, aprendendo a ler diferentes registros escritos, iconográficos, sonoros etc ; VII – Valorizar o patrimônio sociocultural e respeitar a diversidade reconhecendo -a como direito dos povos e indivíduos e como um elemento de fortalecimento da democracia.
Assim, a organização dos conteúdos deve proporcionar o desenvolvimento de uma postura \textbf{crítica} do educando, diante dos elementos que constituem os diferentes contextos tratados pela história.
Nos primeiros anos, temas como a história pessoal da criança, família, tradições e cultura da \textbf{localidade}, permitem estabelecer \textbf{inúmeras} relações, possibilitando às crianças ampliar a \textbf{compreensão} de sua própria história, de suas formas de viver e de se relacionar na sociedade.
Nos anos \textbf{finais} deve -se ampliar o campo de estudos para experiências históricas, no tempo e no espaço.
É preciso \textbf{sistematizar} e programar um ensino em que o \textbf{aluno} parta das experiências vividas e estabeleçam um diálogo com o conhecimento histórico produzido.
Nesse período a referência aos fatos históricos vai se tornando cada vez mais abstrata, sem \textbf{perder} de \textbf{vista} seu vínculo concreto com o espaço-tempo.
O \textbf{fundamento} geral dessa fase aponta para a construção paulatina de macroestruturas históricas permanecendo a necessidade da ordenação, do registro e do desenvolvimento da \textbf{capacidade} argumentativa.
Escrever, desenhar, falar ou representar em busca do desvendamento de causas e consequências, tendo como motivadores os fatos históricos, deverá proporcionar os elementos básicos para o aprimoramento das noções de tempo e suas relatividades.
Dentre as \textbf{competências} e habilidades que devem ser priorizadas para o sexto e sétimo ano estão : àquelas relacionadas à \textbf{capacidade} de compreender as interpelações no tempo e espaço ; As relativas à \textbf{capacidade} de valorizar as peculiaridades culturais dos agrupamentos humanos nas suas diversas temporalidades ; Àquelas relacionadas ao processo de apropriação das relações de dependência e exploração \textbf{econômica} ; As vinculadas ao entendimento dos processos de retroalimentação entre os movimentos históricos de caráter local, nacional e global em diversos tempos e espaços ; As ligadas ao conhecimento das lutas sociais como via legitima de reivindicação na conquista por direitos no Brasil e no Mundo ; Àquelas que \textbf{dizem} respeito ao conhecimento das organizações sociopolíticas, bem como as relações de poder em diversos tempos e espaços ; As que se encontram associadas à \textbf{compreensão} das dinâmicas do processo de colonização da América, África, Ásia e Europa ; As que \textbf{dizem} respeito ao reconhecimento dos espaços participativos como canal \textbf{legítimo} no processo de aprendizagem ; e as que se encontram relacionadas à valorização da liberdade, da vida, da solidariedade, da alteridade, da diversidade e da igualdade como princípios e direitos de todo ser humano.
Nos anos \textbf{finais} o que se deve buscar é a inversão da estrutura \textbf{discursiva}, isto é, procurar relações com o passado a partir de fatos presentes a fim de se ampliar a \textbf{capacidade} de construção \textbf{conceitual} e de ordenação do pensamento do educando, rompendo, assim, com a tradicional linearidade da estrutura \textbf{discursiva}, pautada em relações de causa e efeito lineares.
Inserida nesse processo está à preocupação mais direcionada em formar a identidade e a cidadania social.
A primeira fundada no vínculo entre o passado do grupo de convívio e a história da população brasileira.
A última requer a \textbf{compreensão} de que formar o cidadão envolve prepará -lo para a participação social e política, para o exercício de direitos e deveres políticos, civis e sociais, e para o desenvolvimento de hábitos e valores compatíveis com tal exercício ( FONSECA, 2012 ; OLIVEIRA, 2006, p. 78 ).
Para a consecução dessas finalidades o ensino de história deve considerar três aspectos fundamentais : A relação entre o particular e o geral ; Noções de diferenças e semelhanças, relativas à \textbf{compreensão} do eu ( nós ) e a percepção do outro ( alteridade ) ; Noções de continuidade e permanência ( BRASIL, 1998 ).
Entretanto, ao se ensinar história com o \textbf{foco} na construção de identidades é preciso problematizar o caráter da identidade que se quer formar.
Segundo Silva e Fonseca ( 2012 ) discutir o \textbf{problema} das identidades culturais no contexto das \textbf{pluralidades} requer analisar a rede de significações que a problemática possui em diferentes lugares.
No caso do Brasil significa questionar as implicações de pensar a identidade e a \textbf{pluralidade} cultural, pelo viés da diversidade, descolada da desigualdade social e da discriminação.
Falar da diversidade cultural do Brasil significa levar em conta a origem das famílias e reconhecer as diferenças entre os referenciais culturais.
Significa também reconhecer que encontramos indivíduos que não são iguais que têm especificidades de gênero, raça/etnia, religião, orientações sexuais, valores e outras diferenças definidas a partir de suas histórias pessoais.
Todavia, o respeito à diferença não pode significar o mascaramento ou a omissão perante as profundas desigualdades sociais e \textbf{econômicas} existentes no Brasil.

% matched lemmas: abordagem, aluno, análise, atual, aula, caminho, capacidade, categoria, competência, complexo, compreensão, conceitual, convir, crítico, debate, demanda, depender, desvendar, discursivo, dizer, econômico, epistemologia, esperar, exigir, final, foco, fundamento, identitária, inerente, inúmero, legítimo, letramento, localidade, maneira, método, parâmetro, perder, pluralidade, problema, processar, razão, relevante, sistematizar, teórico, vista
\end{textsample}
